\documentclass[11pt,a4paper]{article}
\usepackage[margin=1in]{geometry}
\usepackage{amsmath,amssymb,amsthm,mathtools,booktabs,microtype,xcolor}
\usepackage[numbers,sort&compress]{natbib}
\usepackage[colorlinks=true,linkcolor=blue!50!black,citecolor=blue!50!black,urlcolor=blue!50!black]{hyperref}
\newtheorem{theorem}{Theorem}[section]
\newtheorem{proposition}[theorem]{Proposition}
\newtheorem{lemma}[theorem]{Lemma}
\newtheorem{corollary}[theorem]{Corollary}
\theoremstyle{remark}\newtheorem{remark}[theorem]{Remark}
\newcommand{\E}{\mathbb E}
\newcommand{\Prb}{\mathbb P}
\newcommand{\Tr}{\operatorname{Tr}}
\newcommand{\R}{\mathbb R}
\newcommand{\Cum}{\operatorname{Cum}}
\title{Thermal Point-Process Statistics:\\Exponential Bounds, Reconstruction, and Their Limits}
\author{}
\date{29 October 2026}
\hypersetup{pdfauthor={},pdftitle={Thermal Point-Process Statistics: Exponential Bounds, Reconstruction, and Their Limits}}
\begin{document}
\maketitle
\begin{abstract}
We separate three questions about a random spectrum: determination of a single thermal statistic by its moments, reconstruction of a point-process law from joint temperature data, and purity of a state representing that law.
For nonnegative bounded weights and factorial-correlation domination, we give a Poisson-type exponential-moment bound and an explicit concentration inequality.
A lognormal Cox construction shows that a nonnegative spectrum with a fixed ground level can nevertheless have a moment-indeterminate heat statistic when correlation domination is absent.
Two deterministic two-level spectra show that a determinate single-clock law need not identify the spectrum.
For determinantal statistics we state the ordered-clock trace inequality and prove that superlinear logarithmic heat-trace growth is sufficient for fixed-order single-channel dominance.
An exactly pure projection process with linear logarithmic growth disproves the converse.
Finally, all mixed positive-temperature moments determine the point-process law in a class with factorial domination and finite weighted transforms, but this does not imply a pure state or a one-dimensional GNS representation.
A conditional semigroup-density lemma isolates what absolute continuity can establish without an unproved stochastic-Airy tail estimate.
The ingredients are classical generating-functional, trace, and moment methods; the results provide explicit assumptions and counterexamples for their use together.
\end{abstract}

\section{Three identification problems}
Let $\xi$ be a random locally finite spectrum and let
\begin{equation}\label{eq:heat}
 Z(\beta)=\int e^{-\beta E}\,\xi(dE),\qquad\beta>0,
\end{equation}
whenever the integral is finite.
The moments of $Z(\beta)$ concern a probability law on $[0,\infty)$.
That law is a coarse image of the random spectrum, while a state on an algebra of observables is a further mathematical object with its own purity question.
Conflating these levels can turn a valid scalar moment calculation into an unsupported spectral or physical conclusion.
The distinction is particularly relevant when gravitational amplitudes are interpreted as ensemble moments~\citep{JanssenMirbabayiZograf2021}.

This paper gives a sequence of positive results and explicit boundary examples.
Sections~\ref{sec:bounded} and~\ref{sec:counter} concern bounded statistics and their assumptions.
Section~\ref{sec:channels} treats determinantal trace channels and their large-clock behavior.
Section~\ref{sec:reconstruct} identifies the stronger joint data sufficient for point-process reconstruction.
Section~\ref{sec:density} explains the limited role of a conditional-density reduction.
The proofs are included so that the scope of each conclusion can be inspected without relying on a proposed gravity completion.

We use standard point-process and determinantal machinery~\citep{Soshnikov2000,Lyons2003}.
No new general theory of moment problems or determinantal processes is claimed.
The useful content is an explicit collection of compatible sufficient conditions, quantitative bounds, and counterexamples to tempting converse implications.

\section{Bounded statistics under correlation domination}\label{sec:bounded}
Let $X$ be a locally compact second-countable space and $\xi$ a simple locally finite point process.
Its factorial correlation measures $\rho_k$ are characterized by
\begin{equation}
 \E\sum_{x_1,\ldots,x_k\in\xi}^{\ne}F(x_1,\ldots,x_k)
 =\int F\,d\rho_k
\end{equation}
for nonnegative measurable $F$; the sum is over ordered distinct points.
Assume the measure inequality
\begin{equation}\label{eq:domination}
 \rho_k\le\rho_1^{\otimes k},\qquad k\ge1.
\end{equation}
For a Hermitian determinantal process with positive-contraction kernel, the Gram determinant inequality gives~\eqref{eq:domination}~\citep{Soshnikov2000,Lyons2003}.
Merely having a formal non-Hermitian determinantal representation is not the same hypothesis.

\begin{theorem}[Exponential-moment bound]\label{thm:mgf}
Suppose~\eqref{eq:domination} holds, $0\le f\le b<\infty$, and
\begin{equation}
 m=\int_X f\,d\rho_1<\infty.
\end{equation}
Then $Z_f=\int f\,d\xi$ is finite almost surely.
For $b>0$ and $t\ge0$,
\begin{equation}\label{eq:mgf}
 \E e^{tZ_f}\le
 \exp\left[\frac mb(e^{tb}-1)\right].
\end{equation}
Its moment generating function extends to an entire function of a complex argument, and its scalar law is determined by its moments.
If $b=0$ or $m=0$, then $Z_f=0$ almost surely.
\end{theorem}
\begin{proof}
Since $\E Z_f=m$, finiteness follows immediately.
First restrict $f$ to a compact window $B$.
Write $g=e^{tf}-1\ge0$ on that window.
The finite-configuration product expansion and Tonelli's theorem give
\begin{align}
 \E e^{tZ_{f1_B}}
 &=\sum_{k=0}^{\infty}\frac1{k!}\int_{B^k}\prod_{i=1}^k g(x_i)\,d\rho_k\\
 &\le\exp\left(\int_Bg\,d\rho_1\right).
\end{align}
Convexity on $[0,b]$ gives
\begin{equation}
 e^{tf}-1\le\frac f b(e^{tb}-1).
\end{equation}
An increasing compact exhaustion and monotone convergence prove~\eqref{eq:mgf}.
For each $R>0$, $|e^{zZ_f}|\le e^{RZ_f}$ on $|z|\le R$ and the right-hand side is integrable.
Differentiation under the expectation therefore defines an entire function, whose Taylor coefficients are the moments.
For any competing probability law on $[0,\infty)$ with the same moments, Tonelli's theorem gives the same finite series $\sum_{n\ge0}t^n\E Z_f^n/n!$ for every $t\ge0$.
Thus both laws have equal moment generating functions near zero and hence equal Laplace transforms; uniqueness of the latter identifies the laws.
\end{proof}

\begin{corollary}[Explicit upper tail]\label{cor:tail}
For $m,b>0$ and $x>0$,
\begin{equation}\label{eq:tail}
 \Prb(Z_f\ge m+x)
 \le\exp\left[-\frac{m+x}{b}\log\left(1+\frac xm\right)+\frac xb\right].
\end{equation}
\end{corollary}
\begin{proof}
Apply Markov's inequality to $e^{tZ_f}$ and~\eqref{eq:mgf}.
The minimizing value is $t=b^{-1}\log(1+x/m)$.
\end{proof}

For the heat weight in~\eqref{eq:heat}, a deterministic lower spectral bound $E\ge E_{\min}$ supplies $b=e^{-\beta E_{\min}}$.
The lower edge need not be zero.
Correlation domination, a uniform finite lower edge, and a finite heat mean suffice for Theorem~\ref{thm:mgf}.
A fixed finite protected degeneracy adds a deterministic constant, shifting the statistic without affecting moment determinacy.
These are hypotheses on a completion, not a construction of one.

\section{Two sharp limits on the interpretation}\label{sec:counter}
\subsection{A positive spectrum with an indeterminate heat law}
The lower-edge hypothesis alone cannot replace correlation control.
We use the classical lognormal moment-indeterminacy mechanism~\citep{Heyde1963} in a simple Cox process.

Let $X$ be Gaussian with mean zero and variance four and set $\Lambda=e^X$.
For $|\epsilon|\le1$, define a density on $(0,\infty)$ by
\begin{equation}\label{eq:stieltjes}
 p_\epsilon(\lambda)=p_0(\lambda)
 \left[1+\epsilon\sin\left(\frac\pi2\log\lambda\right)\right],
\end{equation}
where $p_0$ is the variance-four lognormal density.
The densities are normalized and nonnegative.
For every integer $n\ge0$,
\begin{equation}\label{eq:sinmom}
 \E[e^{nX}\sin(\pi X/2)]
 =e^{2n^2-\pi^2/2}\sin(2\pi n)=0.
\end{equation}
They are distinct probability laws with the same integer moments.

Conditionally on $\Lambda$, sample a Poisson process with intensity $\Lambda\,dE$ on $[0,1]$ and append one fixed level at $E=0$.
Every realization is a finite nonnegative spectrum; almost surely the Poisson part has no point at the appended level.
At $\beta=1$, let
\begin{equation}
 Z=1+\sum_i e^{-E_i}.
\end{equation}
\begin{proposition}\label{prop:cox}
All integer moments of $Z$ are finite and independent of $\epsilon$, but its probability law depends on $\epsilon$.
Consequently each of these heat laws is moment-indeterminate.
\end{proposition}
\begin{proof}
Conditioned moments of the Poisson linear statistic are finite polynomials in $\Lambda$, as follows either by partitioning coincident indices or by differentiating its conditional generating functional.
Equation~\eqref{eq:sinmom} therefore makes all moments of $Z$ equal.
For $s>0$, its Laplace transform is
\begin{equation}\label{eq:coxlap}
 \E_\epsilon e^{-sZ}=e^{-s}L_\epsilon(c(s)),\quad
 c(s)=\int_0^1(1-e^{-s e^{-E}})\,dE,
\end{equation}
where $L_\epsilon$ is the Laplace transform of $\Lambda$.
The function $c$ is continuous, strictly increasing, and maps $(0,\infty)$ onto $(0,1)$.
If two heat laws agreed, their intensity Laplace transforms would agree on $(0,1)$.
Analyticity on the positive half-plane and Laplace uniqueness would identify the distinct densities~\eqref{eq:stieltjes}, a contradiction.
\end{proof}

The missing hypothesis is visible on the diffuse Poisson part:
\begin{equation}
 \frac{\rho_2}{\rho_1\otimes\rho_1}
 =\frac{\E\Lambda^2}{(\E\Lambda)^2}=e^4>1.
\end{equation}
The fixed ground level does not restore factorial domination.
This example therefore does not contradict Theorem~\ref{thm:mgf}.
It shows why a protected lower edge, without a correlation hypothesis, cannot establish the conclusion.

\begin{corollary}\label{cor:coxjoint}
The processes indexed by $\epsilon$ have the same complete mixed-temperature moment hierarchy, despite their distinct heat and point-process laws.
\end{corollary}
\begin{proof}
Conditionally on $\Lambda$, the coincidence-partition formula for any product $\prod_j Z(\beta_j)$ is a finite polynomial in $\Lambda$.
A block contributes the finite coefficient $\int_0^1e^{-(\sum_j\beta_j)E}\,dE$, and the appended ground level only adds constants.
All integer intensity moments agree by~\eqref{eq:sinmom}, whereas Proposition~\ref{prop:cox} already separates the heat laws.
\end{proof}

Thus even all mixed clocks do not identify an unrestricted point-process law.
The reconstruction theorem below uses correlation domination as an essential class assumption, not as incidental regularity.

For an illustration, the portable quadrature gives $\E e^{-Z}\simeq0.1922416392$ and $0.2081849296$ for $\epsilon=1/2$ and $-1/2$, respectively.
These decimals are not needed for the distinct-law proof and are not interval-certified bounds.

\subsection{A determinate clock that does not determine the spectrum}
Consider deterministic two-level Hamiltonians whose Boltzmann weights at $\beta=1$ are
\begin{equation}
 (e^{-E_1},e^{-E_2})=(1/4,3/4),\qquad
 (e^{-F_1},e^{-F_2})=(1/3,2/3).
\end{equation}
Both spectra are positive and have the same dimension.
Both give $Z(1)=1$, so every single-clock moment is exactly one and its law is determinate.
At $\beta=2$, however,
\begin{equation}
 Z_E(2)=\frac58,\qquad Z_F(2)=\frac59.
\end{equation}
Thus even perfect knowledge of a determinate scalar law at one temperature does not reconstruct the spectrum.
No stochastic or infinite-dimensional subtlety is responsible for this failure.

\section{Determinantal channels and large-clock behavior}\label{sec:channels}
\subsection{Cumulant expansion and the two-clock inequality}
For a finite Hermitian positive-contraction kernel $K$ and a nonnegative diagonal weight $f$, the determinantal generating identity is
\begin{equation}
 \E e^{u\int f\,d\xi}=\det(I+(e^{uf}-I)K).
\end{equation}
Expanding the logarithm gives the standard cumulant formula~\citep{Soshnikov2000}
\begin{equation}\label{eq:channels}
 \Cum_k=\sum_{m=1}^k\frac{(-1)^{m+1}}m
 \sum_{\substack{k_1+\cdots+k_m=k\\k_i\ge1}}
 \frac{k!}{k_1!\cdots k_m!}
 \Tr(f^{k_1}K\cdots f^{k_m}K).
\end{equation}
The inner sum is over ordered compositions, so no additional cyclic factorial is needed.
In particular,
\begin{align}
 \Cum_2&=\Tr(f^2K)-\Tr(fKfK),\\
 \Cum_3&=\Tr(f^3K)-3\Tr(f^2KfK)+2\Tr((fK)^3),\\
 \Cum_4&=\Tr(f^4K)-4\Tr(f^3KfK)-3\Tr(f^2Kf^2K)\nonumber\\
 &\hspace{8mm}+12\Tr(f^2K(fK)^2)-6\Tr((fK)^4).
\end{align}
The order-$k$ factorial correlation is a principal determinant; a connected cumulant is the channel sum~\eqref{eq:channels}, not that single determinant.

These formulas also apply to a locally trace-class DPP with unbounded $f$ at each fixed order for which $K^{1/2}f^jK^{1/2}$ is trace class, $1\le j\le k$.
To see this, increase bounded compactly supported weights $f_N$ to $f$.
The factorial expansion bounds the $k$th moments by finite sums of products of $\Tr(f^jK)$, so the moments converge.
The positive operators $K^{1/2}f_N^jK^{1/2}$ converge in trace norm, and telescoping the finite trace products passes~\eqref{eq:channels} to the limit.
This finite-order argument does not assume existence of a positive-$u$ moment generating function for an unbounded heat weight.

\begin{proposition}[Similarly ordered clocks]\label{prop:ordered}
Let $0\le K\le I$ be a finite Hermitian matrix and let $f,g$ be nonnegative diagonal matrices satisfying $(f_i-f_j)(g_i-g_j)\ge0$ for all $i,j$.
Then
\begin{equation}\label{eq:ordered}
 0\le\Tr(fKgK)\le\Tr(fgK).
\end{equation}
\end{proposition}
\begin{proof}
The lower bound is $\Tr(fKgK)=\|g^{1/2}Kf^{1/2}\|_F^2$.
For the upper bound, Hermiticity and symmetrization give
\begin{align}
 \Tr(fgK)-\Tr(fKgK)
 &=\sum_i f_ig_i(K-K^2)_{ii}\nonumber\\
 &\quad+\frac12\sum_{i,j}|K_{ij}|^2(f_i-f_j)(g_i-g_j)\ge0.
\end{align}
\end{proof}

Two positive-temperature clocks increasing in the same spectral coordinate are similarly ordered.
If the denominator is positive, the normalized two-clock cumulant is therefore
\begin{equation}
 0\le 1-\frac{\Tr(fKgK)}{\Tr(fgK)}\le1.
\end{equation}
The ordering condition cannot be dropped.
With $K$ the $2\times2$ all-ones matrix divided by two, $f=\operatorname{diag}(10,1)$, and $g=\operatorname{diag}(1,10)$, the two traces in~\eqref{eq:ordered} are $121/4$ and $10$, respectively.

\subsection{A sufficient growth criterion}
Use a common spectral coordinate $a$ and let $f_\beta=e^{\beta a}$.
Write
\begin{equation}
 \omega(\beta)=\Tr(f_\beta K),\qquad
 \Phi_k(\beta)=\frac{\Cum_k(Z(\beta))}{\omega(k\beta)}.
\end{equation}
This uses $a=-E$ relative to~\eqref{eq:heat}.
The $m=1$ term in~\eqref{eq:channels} is precisely $\omega(k\beta)$.

\begin{theorem}[Superlinear growth is sufficient]\label{thm:condensation}
Suppose all positive-temperature weighted traces above are finite and nonzero, and
\begin{equation}\label{eq:growth}
 \log\omega(\beta)=c\beta^p+o(\beta^p),\qquad c>0,\quad p>1.
\end{equation}
Then $\Phi_k(\beta)\to1$ for every fixed $k\ge2$ as $\beta\to\infty$.
\end{theorem}
\begin{proof}
For a nontrivial composition $k=k_1+\cdots+k_m$, $m\ge2$, let
\begin{equation}
 B_i=K^{1/2}f_\beta^{k_i}K^{1/2}\ge0.
\end{equation}
Schatten H\"older and $\|B_i\|_m\le\|B_i\|_1$ imply
\begin{equation}
 |\Tr(B_1\cdots B_m)|\le\prod_i\|B_i\|_m
 \le\prod_i\omega(k_i\beta).
\end{equation}
The logarithm of this upper bound relative to $\omega(k\beta)$ is
\begin{equation}
 -c\beta^p\left(k^p-\sum_i k_i^p\right)+o(\beta^p).
\end{equation}
Strict superadditivity for $p>1$ makes it tend to minus infinity.
There are finitely many channels at fixed $k$, proving the result.
\end{proof}

\subsection{The converse fails even for a projection kernel}
A non-suppression conclusion cannot be drawn from an upper bound when $p\le1$.
The following pure determinantal example proves that this is a real limitation.

Let $d_n=2^{-n}$, $x_n=1-1/n$, and $y_n=-n$, for $n\ge1$.
On the countable paired-label space, define
\begin{equation}
 v_n=\begin{pmatrix}\sqrt{d_n}\\\sqrt{1-d_n}\end{pmatrix},\qquad
 K=\bigoplus_{n\ge1}v_nv_n^*,\qquad
 f_\beta|_n=\operatorname{diag}(e^{\beta x_n},e^{-\beta n}).
\end{equation}
This is an orthogonal projection DPP: independently in each block exactly one point is selected, at $x_n$ with probability $d_n$ and at $y_n$ with probability $1-d_n$.
All weighted traces at positive temperature are finite.
The selected $x_n$ near the accumulation point $1$ have finite expected number $\sum_n d_n$, hence almost surely finite number, while the $y_n$ escape to minus infinity.
Thus the displayed real spectral configuration is locally finite almost surely.

\begin{proposition}\label{prop:p1}
This projection process satisfies
\begin{equation}
 \lim_{\beta\to\infty}\frac{\log\omega(\beta)}\beta=1,
 \qquad\lim_{\beta\to\infty}\Phi_k(\beta)=1\quad(k\ge2\text{ fixed}).
\end{equation}
\end{proposition}
\begin{proof}
Here
\begin{equation}
 \omega(\beta)=\sum_n\left[d_ne^{\beta x_n}+(1-d_n)e^{-\beta n}\right].
\end{equation}
It is bounded above by $e^\beta+\sum_ne^{-\beta n}$ and below by $d_Me^{\beta(1-1/M)}$ for every fixed $M$.
Taking logarithms, then $\beta\to\infty$ and $M\to\infty$, gives the growth exponent $p=1$.

For a composition of total order $k$, the rank-one block identity gives the channel
\begin{equation}\label{eq:blockchannel}
 \sum_n\prod_{i=1}^m\left[d_ne^{k_i\beta x_n}+(1-d_n)e^{-k_i\beta n}\right].
\end{equation}
Let $D_k=\sum_n d_ne^{k\beta x_n}$.
The all-$x$ term divided by $D_k$ is the weighted average of $d_n^{m-1}$ under weights $d_ne^{k\beta x_n}$.
For any finite initial set the normalized weight tends to zero: a single larger index has strictly greater $x_n$ and eventually dominates that finite set.
Since $d_n^{m-1}\to0$, splitting the average at a large fixed cutoff proves convergence to zero for $m\ge2$.

Any other term in~\eqref{eq:blockchannel} has total $y$ order $k_Y\ge1$ and $x$ order $k_X=k-k_Y$.
Dropping the probability factors bounds it by
\begin{equation}
 e^{k_X\beta}\sum_{n\ge1}e^{-k_Y\beta n}
 \le\frac{e^{(k-2)\beta}}{1-e^{-\beta}}.
\end{equation}
It is negligible compared with $\omega(k\beta)=e^{(k-o(1))\beta}$.
All nontrivial channels vanish after normalization, which proves the claim.
\end{proof}

Thus $p>1$ is sufficient and is not necessary, even with exact projection purity.
The example is infinite rank with finite weighted traces; a finite fixed spectrum cannot realize the superlinear growth premise~\eqref{eq:growth} in the first place.
Neither this channel limit nor rapid scalar moment growth establishes moment indeterminacy.
In fact $f_\beta\le e^\beta$ in this projection example, and its finite mean and determinantal correlation domination put each fixed-$\beta$ law under Theorem~\ref{thm:mgf}.
Its scalar laws are therefore moment-determinate while its normalized fixed-order cumulants still tend to one.

\section{What all joint clocks determine}\label{sec:reconstruct}
We now specify the stronger data that recover a spectral process.
Use the coordinate $a=-E$ and
\begin{equation}
 Z(\beta)=\int e^{\beta a}\,\xi(da).
\end{equation}
\begin{theorem}[Joint-temperature reconstruction]\label{thm:reconstruction}
Let $\xi$ be a simple locally finite point process on $\R$ whose first intensity is locally finite and whose factorial correlations satisfy~\eqref{eq:domination}.
Assume
\begin{equation}\label{eq:lapfinite}
 \int_\R e^{\beta a}\,\rho_1(da)<\infty
 \quad\text{for every }\beta>0.
\end{equation}
Within this class, the complete collection
\begin{equation}\label{eq:jointdata}
 \left\{\E\prod_{i=1}^k Z(\beta_i):k\ge1,\ \beta_i>0\right\}
\end{equation}
determines the point-process law uniquely.
\end{theorem}
\begin{proof}
Partition the indices $1,\ldots,k$ according to which factors use the same point of $\xi$.
For a partition $\pi$ with blocks $B$, this gives
\begin{equation}\label{eq:partition}
 \E\prod_{i=1}^kZ(\beta_i)
 =\sum_{\pi}\int\prod_{B\in\pi}
 e^{(\sum_{i\in B}\beta_i)a_B}\,d\rho_{|\pi|}((a_B)_{B\in\pi}).
\end{equation}
All terms are finite by~\eqref{eq:domination} and~\eqref{eq:lapfinite}.
The partition into singletons is the $k$-variable Laplace transform of $\rho_k$.
All other terms use lower-order correlation measures.
Induction therefore recovers
\begin{equation}
 L_k(\beta_1,\ldots,\beta_k)
 =\int e^{\sum_i\beta_i a_i}\,d\rho_k
\end{equation}
from~\eqref{eq:jointdata}.

Fix $\beta_0>0$ and tilt $\rho_k$ to the finite measure
\begin{equation}
 d\mu_k=e^{\beta_0\sum_i a_i}\,d\rho_k.
\end{equation}
Its moment generating function is known and finite on the open box $|h_i|<\beta_0$, by evaluating $L_k$ at $\beta_0+h_i$.
Multivariate Laplace uniqueness determines $\mu_k$, and untilting determines $\rho_k$.
This step needs no one-sided support assumption.

It remains to identify the process from the recovered correlations.
For bounded $B$, set $\alpha_B=\rho_1(B)<\infty$.
Then $\rho_k(B^k)\le\alpha_B^k$.
For $|h|\le1$ supported in $B$, the factorial generating-functional expansion is absolutely convergent, bounded in absolute sum by $e^{\alpha_B}$:
\begin{equation}
 \E\prod_{a\in\xi}(1+h(a))
 =\sum_{k=0}^{\infty}\frac1{k!}\int\prod_{i=1}^kh(a_i)\,d\rho_k.
\end{equation}
Taking $h=e^{-v}-1$ for nonnegative $v$ supported in $B$ gives the local Laplace functional.
The collection of these functionals determines the locally finite point-process law.
\end{proof}

The proof uses the complete mixed hierarchy, not independent knowledge of the marginal laws at each temperature.
An unbounded heat weight can fail to have positive exponential moments even though every finite mixed moment in~\eqref{eq:jointdata} exists.
The proof does not apply Theorem~\ref{thm:mgf} outside its bounded-weight hypothesis.

Airy spectral coordinates have polynomial intensity growth toward $a\to-\infty$ and rapid decay toward $a\to+\infty$.
The corresponding heat convention is $e^{+\beta a}$.
If instead physical energies $E=-a$ are used, both tail directions and the exponential sign must be reversed together.
Theorem~\ref{thm:reconstruction} is a class-restricted law-identification statement; it does not determine which physical completion supplies its hypotheses.

\subsection{Unique law is not a pure state}
Take the probability measure $\mu=(\delta_0+\delta_1)/2$ and its polynomial state $\omega(p)=\int p\,d\mu$.
Its moments uniquely determine $\mu$: any positive representing measure must obey
\begin{equation}
 \int(x^2-x)^2\,d\mu=0,
\end{equation}
so it is supported on $\{0,1\}$, and its first moment fixes both weights.
But its GNS Hilbert space is $L^2(\mu)$, of dimension two, and
\begin{equation}
 \omega=\tfrac12\operatorname{ev}_0+\tfrac12\operatorname{ev}_1
\end{equation}
is not pure.
Thus a unique representing distribution is not a single spectral point.
A nondegenerate determinate Gaussian provides the same distinction with an infinite-dimensional $L^2$ representation.
Joint-temperature reconstruction therefore does not imply a one-dimensional representation, a pure ensemble state, or a conclusion about unitary time evolution.

\section{A conditional density method and its exact ceiling}\label{sec:density}
Scalar moment questions with an unbounded heat weight may call for information about a density rather than exponential moments.
The following elementary semigroup method isolates a useful part of that problem.

\begin{lemma}[One Gaussian direction]\label{lem:semigroup}
Let $H_\perp$ be lower-bounded self-adjoint on $L^2(0,\infty)$.
Let $u\in L^\infty$ be real, compactly supported, nonpositive, and nonzero on a set of positive measure, and let $c,\beta>0$.
Assume $H_\tau=H_\perp+c\tau M_u$ has trace-class heat semigroup for all real $\tau$, its diagonal heat kernel is strictly positive almost everywhere on the support of $-u$, and bounded-perturbation differentiation is valid in trace norm.
Then $g(\tau)=\Tr e^{-\beta H_\tau}$ is continuously differentiable and
\begin{equation}\label{eq:derivative}
 0<g'(\tau)=-\beta c\Tr(M_u e^{-\beta H_\tau})
 \le\beta c\|u\|_\infty g(\tau).
\end{equation}
If $\tau$ is an independent standard Gaussian coordinate, conditionally on the complementary randomness the law of $g(\tau)$ has density on its image
\begin{equation}\label{eq:conditional}
 p(z\mid\perp)=\frac{\varphi(\tau_z)}{g'(\tau_z)}
 \ge\frac{\varphi(\tau_z)}{\beta c\|u\|_\infty z},\qquad g(\tau_z)=z.
\end{equation}
\end{lemma}
\begin{proof}
Duhamel differentiation gives
\begin{equation}
 \frac{d}{d\tau}e^{-\beta H_\tau}
 =-c\int_0^\beta e^{-(\beta-r)H_\tau}M_u e^{-rH_\tau}\,dr.
\end{equation}
Trace cyclicity makes the integrand independent of $r$.
Positivity of the diagonal and $0\le-u\le\|u\|_\infty$ give~\eqref{eq:derivative}.
Strict monotonicity permits the one-dimensional change-of-variable formula, proving~\eqref{eq:conditional}.
Mixtures of these conditional densities remain absolutely continuous.
\end{proof}

The lemma differentiates the heat trace as a whole and avoids level-by-level eigenfunction derivatives.
It is useful only with the displayed pathwise and Gaussian-disintegration hypotheses.
The stochastic Airy operator is an established representation of soft-edge spectra~\citep{RamirezRiderVirag2011}; applying Lemma~\ref{lem:semigroup} to a chosen realization still requires checking those hypotheses.

More substantially,~\eqref{eq:conditional} is not a uniform pointwise tail lower bound.
At a prescribed large or small $z$, one must control the Gaussian inverse coordinate $\tau_z$ on a quantitatively positive-mass set of complementary slices that all reach that $z$.
A positive probability of landing somewhere in a broad band does not supply this common incidence estimate.
An upper derivative bound cannot supply a lower bound on the image length traversed during a fixed parameter interval.
Spectral control also requires correctly directed bracketing: an inserted internal Dirichlet condition raises the quadratic form and cannot be used as an upper bound on the original heat trace.

If suitable two-ended pointwise density bounds and compact-interval lower bounds were proved, a classical logarithmic Krein criterion could imply moment indeterminacy.
Neither Lemma~\ref{lem:semigroup}, the growth assumption~\eqref{eq:growth}, nor failure of a sufficient Carleman criterion proves those bounds.
We make no unconditional fixed-temperature stochastic-Airy indeterminacy assertion.

\section{Interpretation and reproducibility}
The results distinguish four levels of information.
Correlation domination and a bounded weight yield an entire scalar moment generating function and a quantitative tail bound.
Removing domination permits moment-indeterminate heat laws even on positive spectra with a fixed ground level.
Knowing a determinate single-clock law still need not identify the spectrum, while the full mixed-clock hierarchy identifies a process law in the explicitly bounded correlation class.
That unique law can remain a nontrivial mixture of spectral realizations.

The accompanying portable scripts verify the lognormal signed-moment identity, the two-level example, finite determinantal exponential bounds, low-order cumulant coefficients, and the unordered-clock counterexample.
They record the Cox Laplace quadrature and finite channel-ratio illustrations, clearly separated from the analytic proofs.
They require no archived project path or network access.
The pure-projection channel limit and the full-clock reconstruction are analytic results; finite decimal sequences do not certify them.
No stochastic-Airy simulation or gravitational dynamics is supplied by these tests.

The article is a methods synthesis with explicit assumption tests.
Classical lognormal indeterminacy, determinantal generating functions, and Laplace uniqueness remain credited prior mathematics.
Its statements neither establish a nonperturbative gravity completion nor decide black-hole information loss.

\section*{AI-assisted workflow and unresolved attribution}
The repository operator states:
\begin{quote}
I genuinely do not know how authorship, intellectual credit, or ownership should be attributed for these results.
\end{quote}
The operator reports that earlier work was led mainly by Fable 5, whereas results from July 2026 onward were predominantly produced by GPT models.
DeepSeek and multiple GPT models also participated in cross-review, checking, and manuscript preparation.
This is a reported chronology rather than a complete audited allocation of individual contributions.
A later correction does not reattribute an earlier result.
Human author names and affiliations are intentionally blank, and these scientific claims have not undergone human expert review.
Historical research dates, the current edition date, private upload, and demonstrated public disclosure remain distinct; no settled authorship, ownership, publication-priority, or journal-acceptance claim is made.
\bibliographystyle{plainnat}
\bibliography{references}
\end{document}
